\chapter{Continuous Scale, Fixed Points, Spectra, and Two-Index Limits}
\label{ch:wilsonian-limits}

Given the comparison data of \cref{ch:renormalization}, a renormalization
flow usually carries four further constructions: a continuous scale with its
beta function, fixed points with their linearization, spectra, and limits in
which system size and depth grow together. This chapter asks what each
requires and what survives for the finite networks of the main manuscript.
Over an interval, a continuous beta function is a \(C^1\) family of comparison
maps that discrete endpoints do not determine
(\cref{prop:rg-continuous-beta-underdetermined}). A fixed measure pair of a
rescaled kernel with a unique invariant law has constant action
(\cref{prop:rg-fixed-pair-unique-ergodicity}), and fixedness does not transfer
between the law, action, interaction, configuration, and attention tiers
(\cref{prop:rg-fixed-object-nonimplication}). The positive-unital
essential-spectrum theorem has a hypothesis that no exact arrow satisfies
(\cref{thm:rg-unital-essential-spectrum}), and one exact kernel has different
spectral radii in two topologies (\cref{prop:rg-circle-norm-witness}).
Fixed-depth and maximal-depth limits need not agree
(\cref{prop:rg-noncommuting-limits}). On Gaussian operator hierarchies, dense
scaling leads to a boundary operator ray
(\cref{prop:rg-gaussian-dense-scaling}), and under one primitive
identification Perron--Frobenius attraction fixes the fiber matrix without
selecting it (\cref{conj:grg-fixed-b-attraction}). The open problems of the
chapter state what an infinite-network theory of the agent network would have
to supply.

Beyond the symbols of \cref{ch:renormalization}, this chapter writes \(t\) for
a scale coordinate, \(\mathfrak b^{\mathrm{scale}}\) for the continuous beta,
\((\varrho_*,\mathsf m_*)\) for a fixed measure pair, \(\mathrm T\) for a tier label,
\(\mathcal N\) for spectral counting functions, \(\mathcal W\) for operator
cones, and \(\xi\) for growth exponents of block sums.

\paragraph{Scope.}
The statements this part declines are listed in the scope paragraph of
\cref{ch:renormalization}. The open constructions of this chapter are a
canonical continuous beta (\cref{open:rg-continuous-beta}), two-index limits
and universality (\cref{open:rg-two-index-limits}), and a primitive
hierarchical identification for Gaussian operator hierarchies
(\cref{open:rg-scalarized-attraction}); the continuum theory of laws is
\cref{open:prob-continuum-theory}.

\section{Continuous scale}

\definitionheading{Smooth scale bundle, scale connection, and continuous beta}{def:rg-scale-connection}
A continuous beta function lives on a separately declared tier. Declare an
oriented one-dimensional \(C^1\) scale manifold \(S\) with coordinate \(t\), a
\(C^1\) Banach vector bundle \(\pi_S:\mathscr G\to S\) whose fibers are coupling
spaces, and a scale connection, a covariant derivative
\(\nabla^{\mathrm{scale}}\) on sections of \(\mathscr G\) with the usual
linearity and Leibniz rules. For a \(C^1\) running coupling section
\(g(t)\in\mathscr G_t\), the continuous beta is
\begin{equation}
\mathfrak b^{\mathrm{scale}}(t)=\nabla^{\mathrm{scale}}_{\partial_t}g(t)
\in\mathscr G_t.
\label{eq:rg-covariant-scale-beta}
\end{equation}
In a local \(C^1\) trivialization \(J_t:\mathcal G_\star\to\mathscr G_t\) with
connection form \(A_t\in\mathcal L(\mathcal G_\star)\) and
\(\widetilde g(t)=J_t^{-1}g(t)\),
\begin{equation}
J_t^{-1}\mathfrak b^{\mathrm{scale}}(t)
=\partial_t\widetilde g(t)+A_t\widetilde g(t),
\qquad
A'_t=\mathcal S_tA_t\mathcal S_t^{-1}-(\partial_t\mathcal S_t)\mathcal S_t^{-1}
\ \text{ under }J'_t=J_t\mathcal S_t^{-1}.
\label{eq:rg-scale-connection-form}
\end{equation}
A raw coordinate derivative does not obey this law: on the trivial line bundle
with \(g=1\) parallel in a frame with \(A=0\), the frame change
\(\mathcal S_t=e^{t^2}\) gives raw derivative \(2te^{t^2}\neq0\), while
\(A'_t=-2t\) and the covariant derivative
\(\partial_t(e^{t^2})-2te^{t^2}\) vanishes. The scale connection compares
coupling fibers at different resolutions. It is distinct from the principal
connections of \cref{ch:base-geometry}, which compare statistical fibers along
the base at one resolution, and from the inference-orbit parameter and the
Fisher duration of \cref{ch:relational-inference}.

\propositionheading{Density realization of the continuous beta}{prop:rg-density-beta}
Suppose scale is a continuous semigroup and that, relative to one common
dominating measure, positive densities
\(r_t,p_t\in\operatorname{Dom}(\mathcal A^*)\) are differentiable in a topology
that permits pointwise ratios and the dual pairings below, with
\(p_t=r_t\exp[-H_t]\), \(\dot r_t=\mathcal A^*r_t\), and
\(\dot p_t=\mathcal A^*p_t\). Then, almost everywhere wherever the displayed
ratios are defined,
\begin{equation}
\dot H_t
=-\frac{\mathcal A^*(r_t\exp[-H_t])}{r_t\exp[-H_t]}
+\frac{\mathcal A^*r_t}{r_t},
\label{eq:rg-continuous-beta-functional}
\end{equation}
and the second term vanishes for a stationary reference. If, on a declared
Banach action space with a Schauder basis \(\{\psi_A\}\) and continuous
coordinate functionals \(\{\psi_A^*\}\), or a Hilbert space with a Riesz basis,
the expansion \(H_t=\sum_Ag_A(t)\psi_A\) converges and differentiation and
coefficient extraction can be interchanged, then the component beta functions
are
\begin{equation}
\dot g_A=\left\langle\psi_A^*,
-\frac{\mathcal A^*(r_te^{-H_t})}{r_te^{-H_t}}
+\frac{\mathcal A^*r_t}{r_t}\right\rangle,
\label{eq:rg-component-beta}
\end{equation}
and a truncated basis defines a projection with an explicit residual.
\status{ESTABLISHED}

\noindent\emph{Proof.}
Differentiate \(H_t=-\log(p_t/r_t)\) and substitute the two forward equations;
apply \(\psi_A^*\) to both sides and use the interchange hypothesis. \(\square\)

\propositionheading{Discrete endpoints do not determine a continuous beta}{prop:rg-continuous-beta-underdetermined}
(i) On an interval of \(S\) carrying a \(C^1\) trivialization \(J_t\) whose
connection form \(t\mapsto A_t\) is continuous in operator norm, the equation
\(\partial_t\mathcal S_t=\mathcal S_tA_t\), \(\mathcal S_{t_0}=I\), has a unique
\(C^1\) solution by bounded invertible operators, and in the frame
\(J'_t=J_t\mathcal S_t^{-1}\) the connection form vanishes and
\(\mathfrak b^{\mathrm{scale}}(t)=J'_t\partial_t\bigl(J'^{-1}_tg(t)\bigr)\); over
an interval the scale tier is therefore equivalent to a \(C^1\) family of
comparison isomorphisms. (ii) A scale evolution family
\begin{equation}
\mathsf V(t,t'):\mathscr G_{t'}\longrightarrow\mathscr G_t,
\qquad
\mathsf V(t,t')\circ\mathsf V(t',t'')=\mathsf V(t,t''),
\qquad
\mathsf V(t,t)=I_{\mathscr G_t},
\label{eq:rg-scale-evolution-family}
\end{equation}
the continuous analogue of \eqref{eq:rg-cocycle-composition} and again a
two-parameter law, is not determined by its restriction to integer endpoints,
and neither is \eqref{eq:rg-covariant-scale-beta}: on \(S=\R\) with the trivial
line bundle and zero connection form, the families
\begin{equation}
\mathsf V^{(0)}(t,t')=1,
\qquad
\mathsf V^{(\epsilon)}(t,t')
=\exp\bigl\{\epsilon[\sin(2\pi t)-\sin(2\pi t')]\bigr\}
\label{eq:rg-scale-interpolation-witness}
\end{equation}
both satisfy \eqref{eq:rg-scale-evolution-family} and agree at every pair of
integer endpoints, while their generators are \(0\) and
\(2\pi\epsilon\cos(2\pi t)\). The one-parameter form, in which
\(\mathsf V(t,t')\) depends only on \(t-t'\), is not a consequence of
\eqref{eq:rg-scale-evolution-family}. \status{ESTABLISHED}

\noindent\emph{Proof.}
(i) On a compact subinterval, \(\sup_t\lVert A_t\rVert<\infty\) and
\(\lVert\mathcal SA_t-\mathcal S'A_t\rVert\leq\lVert A_t\rVert
\lVert\mathcal S-\mathcal S'\rVert\), so Picard iteration for the linear equation
in the Banach algebra \(\mathcal L(\mathcal G_\star)\) converges to a unique
\(C^1\) solution, and the solutions on nested compact subintervals agree. The
solution \(\mathcal T_t\) of \(\partial_t\mathcal T_t=-A_t\mathcal T_t\),
\(\mathcal T_{t_0}=I\), satisfies
\(\partial_t(\mathcal S_t\mathcal T_t)=\mathcal S_tA_t\mathcal T_t
-\mathcal S_tA_t\mathcal T_t=0\), so \(\mathcal S_t\mathcal T_t=I\); and
\(\mathcal U_t=\mathcal T_t\mathcal S_t\) solves
\(\partial_t\mathcal U_t=\mathcal U_tA_t-A_t\mathcal U_t\) with
\(\mathcal U_{t_0}=I\), whose unique solution is \(I\), so
\(\mathcal T_t\mathcal S_t=I\). By \eqref{eq:rg-scale-connection-form},
\(A'_t=\mathcal S_tA_t\mathcal S_t^{-1}-\mathcal S_tA_t\mathcal S_t^{-1}=0\).
(ii) Both families have the form \(f(t)/f(t')\) with \(f\) positive, so they
satisfy the composition and identity laws. At integer \(t\) and \(t'\) the sine
terms vanish. Differentiating \(\mathsf V^{(\epsilon)}(t,t')\) in \(t\) at
\(t=t'\) gives \(2\pi\epsilon\cos(2\pi t')\), which is not identically zero for
\(\epsilon\neq0\), so \(\mathsf V^{(\epsilon)}\) is not a function of \(t-t'\).
\(\square\)

The smooth scale base, the bundle, the interpolation, and the connection are
therefore declared input; a discrete blocking sequence supplies the discrete
beta \eqref{eq:rg-interaction-beta} and nothing more.

\openproblemheading{Canonical continuous beta}{open:rg-continuous-beta}
\Cref{def:rg-scale-connection} declares the tier a continuous beta requires,
and \cref{prop:rg-continuous-beta-underdetermined} shows that over an interval
the tier is a \(C^1\) comparison family which a discrete blocking sequence does
not determine. No canonical interpolation of a discrete sequence is known, and
no reported continuous beta has been shown independent of the interpolation
inside a declared class. The problem is settled by a canonical interpolation
with a proof of its uniqueness in a declared class, or by a proof that the
continuous beta is the same for every member of that class.

\section{Fixed points, flows, and scaling operators}
\label{sec:rg-fixed-points}

\definitionheading{Fixed measure pair}{thm:rg-fixed-point-equations}
For a declared family of ratios \(b\) and rescaled kernels \(K_b\) on one
common standard Borel space \(\mathsf Y\) (\cref{def:rg-rescaled-kernel}), a
fixed theory is a measure pair \((\varrho_*,\mathsf m_*)\) with
\(0<\mathsf m_*(\mathsf Y)<\infty\) and
\begin{equation}
\mathfrak R_b(\varrho_*,\mathsf m_*)=(\varrho_*,\mathsf m_*)
\quad\text{for every declared }b.
\label{eq:rg-fixed-measure-pair}
\end{equation}
On the dominated tier, \(\mathsf m_*=e^{-H_*}\varrho_*\), and by uniqueness of
Radon--Nikodym derivatives and \eqref{eq:rg-reference-dependent-action-map} this
is equivalent to
\begin{equation}
\mathfrak R^\varrho_b\varrho_*=\varrho_*,
\qquad
\mathfrak R^H_b[H_*;\varrho_*]=H_*,
\label{eq:rg-fixed-action}
\end{equation}
\(\varrho_*\)-almost everywhere. If the evidence mass is discarded and actions are
taken projectively, the fixed-ray equation is
\begin{equation}
\mathfrak R^H_b[H_*;\varrho_*]=H_*+c_{\mathrm{ray}}(b);
\label{eq:rg-fixed-action-ray}
\end{equation}
with an invariant reference it forces \(c_{\mathrm{ray}}(b)=0\), since \(K_b\)
preserves mass and
\(\mathsf m_*(\mathsf Y)=(\mathsf m_*K_b)(\mathsf Y)=e^{-c_{\mathrm{ray}}(b)}\mathsf m_*(\mathsf Y)\). A
fixed attention law is a fixed rescaled event measure \(\eta_*\); by uniqueness
of disintegration it is fixed exactly when its receiver marginal \(\alpha\) is
fixed and its row \(\beta\) is fixed \(\alpha\)-almost everywhere.

\propositionheading{Fixed measure pairs of a uniquely ergodic kernel have constant action}{prop:rg-fixed-pair-unique-ergodicity}
Suppose that for some declared \(b\) the rescaled kernel \(K_b=C_bI_b\) on the
common standard Borel space has exactly one invariant probability law. If
\((\varrho_*,\mathsf m_*)\) satisfies \eqref{eq:rg-fixed-measure-pair} with
\(0<M_*=\mathsf m_*(\mathsf Y)<\infty\), then \(\mathsf m_*=M_*\varrho_*\), so on the dominated
tier \(H_*=-\log M_*\) is constant \(\varrho_*\)-almost everywhere. On a finite
space the hypothesis holds whenever \(K_b\) is irreducible, in particular
whenever some power of \(K_b\) has all entries positive, and \(K_b\) has all
entries positive whenever every row of \(I_b\) charges every point.
Consequently a fixed theory with nonconstant action requires a rescaled kernel
with at least two invariant probability laws: on a finite space a reducible
rescaled kernel, for example a reducible identification, and otherwise an
infinite space on which uniqueness of the invariant law fails.
\status{ESTABLISHED}

\noindent\emph{Proof.}
The normalized kernel preserves total mass, so \(\mathsf m_*K_b=\mathsf m_*\) makes \(\mathsf m_*/M_*\)
an invariant probability law, and uniqueness gives \(\mathsf m_*/M_*=\varrho_*\). For the
finite case, first let some power \(K_b^p\) have all entries positive. For
probability vectors \(\mu,\nu\), write \(\mu-\nu=s(\mu_1-\nu_1)\) with
\(s=\lVert\mu-\nu\rVert_{\mathrm{TV}}\) and probability vectors
\(\mu_1,\nu_1\); convexity gives
\(\lVert\mu K_b^p-\nu K_b^p\rVert_{\mathrm{TV}}\leq\delta\lVert\mu-\nu\rVert_{\mathrm{TV}}\)
with \(\delta=\max_{y,y'}\lVert K_b^p(y,\cdot)-K_b^p(y',\cdot)\rVert_{\mathrm{TV}}
\leq1-\sum_z\min_yK_b^p(y,z)<1\), so two invariant laws of \(K_b\), which are
invariant for \(K_b^p\), coincide. If \(K_b\) is only irreducible on \(n\)
states, the lazy kernel \((I+K_b)/2\) has the same invariant laws, and
\((I+K_b)^{n-1}=\sum_{k=0}^{n-1}\binom{n-1}{k}K_b^k\) has all entries positive,
because irreducibility gives a positive-probability path of length at most
\(n-1\) between any two states; the previous case applies. Finally,
\((C_bI_b)(y,y')=\sum_zC_b(y,z)I_b(z,y')>0\) when every \(I_b(z,y')>0\), because
each row of \(C_b\) is a probability vector. \(\square\)

At a fixed pair with \(H_*\) in the declared finite-valued action Banach space,
set
\begin{equation}
\mathsf m_*(dy)=e^{-H_*(y)}\varrho_*(dy),\qquad
M_*=\mathsf m_*(\mathsf Y),\qquad
\bar{\mathsf m}_*=M_*^{-1}\mathsf m_*,\qquad
\bar{\mathsf m}^c_*=\bar{\mathsf m}_*K_b,
\label{eq:rg-fixed-normalized-reverse-data}
\end{equation}
fix a standard Borel reverse-kernel version \(K_b^\dagger\) with
\begin{equation}
\bar{\mathsf m}_*(dy)K_b(y,dz)=\bar{\mathsf m}^c_*(dz)K_b^\dagger(z,dy),
\label{eq:rg-fixed-normalized-reverse-kernel}
\end{equation}
and require
\begin{equation}
0<\E_{\varrho_*}\bigl[e^{-H_*}\given Z=z\bigr]<\infty
\quad\varrho_*K_b\text{-almost everywhere}.
\label{eq:rg-linearization-positive-likelihood}
\end{equation}
For bounded perturbations \(\varphi\in L^\infty(\bar{\mathsf m}_*)\),
\cref{thm:rg-bounded-action-calculus}, applied to the pair \((\varrho_*,\mathsf m_*)\) and
the composite kernel \(K_b\), which already includes the rescaling, gives
\(\bar{\mathsf m}^c_*\)-almost everywhere
\begin{equation}
D_H\mathfrak R^H_b\big|_{(H_*,\varrho_*)}[\varphi](z)
=\E_{K_b^\dagger(z,dy)}\varphi(Y),
\qquad
D_H^2\mathfrak R^H_b\big|_{(H_*,\varrho_*)}[\varphi,\varphi](z)
=-\operatorname{Var}_{K_b^\dagger(z,dy)}\varphi(Y).
\label{eq:rg-linearized-action}
\end{equation}

\remarkheading{A declared passive projection scheme}{rem:rg-passive-blocking-computation}
On a homogeneous six-cycle with ratio two and a homogeneous nine-cycle with
ratio three, a declared projected map was iterated: tile one site table and one
pair table onto the cycle, apply the blocking step with attention frozen as
declared data, and read back one coarse site table and one pair table in
anchored M\"obius coordinates. This is a projection scheme on a
finite-dimensional coupling space, not the exact step
\eqref{eq:rg-exact-nonlinear-interaction-map} with comparison isomorphisms, and
no exact-image invariance \eqref{eq:rg-interaction-exact-image-invariance} is
proved for its read-back. At both ratios the iteration converged to fixed
structures of the scheme whose pairwise block vanished to machine precision
(check \texttt{CHK-RG-PASSIVE-BLOCKING}, \cref{tab:base-numerical}); these are
fixed points of the scheme, not fixed objects of the exact step.
\status{NUMERICAL}

\remarkheading{A declared row-regeneration scheme}{rem:rg-regenerated-attention-computation}
On the same instances, a declared row-regeneration rule, which at each level
recomputes attention rows as softmax rows of transported divergences averaged
over the flow, with \(\tau=1\) and uniform receiver occupancies, and folds the
alignment energy into the coarse action, gave fixed structures of the scheme
with nonzero pairwise block at both ratios (check
\texttt{CHK-RG-REGENERATED-ATTENTION}, \cref{tab:base-numerical}). The rule is
a declared approximate channel, not the exact event-law pushforward
\eqref{eq:rg-meta-attention}: its rows are softmax rows of averaged energies,
the form of the constant-row optimum of \cref{prop:obs-attention-constant-row-gap},
which differs from an average of posterior rows by the gap
\eqref{eq:obs-attention-jensen-gap}, and its error relative to the exact
pushforward is not bounded. The passive triviality of
\cref{rem:rg-passive-blocking-computation} is therefore a property of the
frozen-attention scheme on these instances; nothing here shows that the exact
pushforward maintains interaction. \status{NUMERICAL}

\definitionheading{Typed invariant and fixed objects}{def:rg-typed-fixed-objects}
Write \(Y^{\mathrm T}_\ell\) for the object of one declared tier \(\mathrm T\)
and \(F^{\mathrm T}_\ell:Y^{\mathrm T}_\ell\to Y^{\mathrm T}_{\ell+1}\) for its
declared step. The general invariant object of a nonautonomous scale diagram
is an invariant section, a sequence with
\begin{equation}
y_{\ell+1}=F^{\mathrm T}_\ell(y_\ell),
\label{eq:rg-invariant-section}
\end{equation}
which needs no identification. After declared isomorphisms
\(J_\ell:Y^{\mathrm T}_\star\to Y^{\mathrm T}_\ell\), the reference step is
\begin{equation}
\widehat F^{\mathrm T}_\ell
=J_{\ell+1}^{-1}\circ F^{\mathrm T}_\ell\circ J_\ell
:Y^{\mathrm T}_\star\longrightarrow Y^{\mathrm T}_\star,
\label{eq:rg-reference-step}
\end{equation}
and a reference fixed object is a \(y_\star\in Y^{\mathrm T}_\star\) with
\begin{equation}
\widehat F^{\mathrm T}_\ell(y_\star)=y_\star
\qquad\text{for every }\ell,
\label{eq:rg-reference-fixed-object}
\end{equation}
the ordinary fixed-point equation in an autonomous scheme. If the identified
maps are \(p\)-periodic, the base-\(\ell\) reference monodromy is
\begin{equation}
\widehat{\mathcal M}^{\mathrm T}_\ell
=\widehat F^{\mathrm T}_{\ell+p-1}\circ\cdots\circ\widehat F^{\mathrm T}_\ell
:Y^{\mathrm T}_\star\longrightarrow Y^{\mathrm T}_\star,
\label{eq:rg-monodromy-map}
\end{equation}
a reference monodromy fixed object satisfies
\(\widehat{\mathcal M}^{\mathrm T}_\ell(y_{\star,0})=y_{\star,0}\), and its
reference \(p\)-cycle
\(y_{\star,j+1}=\widehat F^{\mathrm T}_{\ell+j}(y_{\star,j})\),
\(y_{\star,p}=y_{\star,0}\), lies in \(Y^{\mathrm T}_\star\), with points that
need not be fixed by any one-step reference map. Without a common object the
equation \(F^{\mathrm T}_\ell(y)=y\) is ill typed, and a fixed ray additionally
requires a declared cone and projectivization. The tiers are the law tier (the
normalized law, the measure pair of \cref{def:rg-measure-pair}, and the event
law \eqref{eq:rg-attention-event-law}, with fixed objects as in
\cref{thm:rg-fixed-point-equations}); the action tier (\(H_\ell\) relative to
\(\varrho_\ell\) and, separately, its additive class, related by
\eqref{eq:rg-fixed-action} and \eqref{eq:rg-fixed-action-ray} and meaningful
only with a reference trajectory by
\cref{prop:rg-action-beta-reference-change}); the interaction tier
(\(\mathsf g_\ell\in\mathcal G_\ell\), step \(\mathcal T_\ell\), reference step
\(\widehat{\mathcal T}_\ell\)); the retained tier
(\(\operatorname{Ran}\mathsf P^{\mathrm{ret}}_\ell\), with exact fixed objects only
under \eqref{eq:rg-beta-exactness-criterion}); the gauge-retained state tier
\eqref{eq:rg-full-state}; and the configuration tier of
\cref{ch:relational-inference}.

\propositionheading{Fixedness does not transfer between tiers}{prop:rg-fixed-object-nonimplication}
None of the following implications holds: a fixed normalized law gives a fixed
reference-dependent action; a fixed action class gives a fixed tracked measure
pair; a fixed retained interaction gives a fixed exact interaction; a fixed law
gives a fixed configuration; a fixed conditional attention row gives a fixed
event law; a monodromy fixed object is a one-step fixed object.
\status{ESTABLISHED}

\noindent\emph{Proof.}
On a two-point set, hold \(\mu=(1/2,1/2)\) fixed while the reference alternates
between \((1/4,3/4)\) and \((3/4,1/4)\); the densities of \(\mu\) are
\((2,2/3)\) and \((2/3,2)\), so the actions are interchanged. The actions \(H\)
and \(H+c\) have the same class, while \(e^{-H}\varrho\) and \(e^{-c}e^{-H}\varrho\)
have evidence masses differing by the factor \(e^{-c}\), which is why
\eqref{eq:rg-fixed-action-ray} is weaker than \eqref{eq:rg-fixed-measure-pair}
until mass preservation forces \(c_{\mathrm{ray}}(b)=0\). The witness of
\cref{prop:rg-retained-beta-residual} has zero retained beta at every \((x,0)\)
and exact image \((x,x)\) for \(x\neq0\). On the configuration manifold
\(S^1\) with the antipodal step \(x\mapsto-x\) and a constant
configuration-to-law extraction, the law is fixed while the configuration step
has no fixed point. Holding \(\beta\) fixed while the receiver occupancy
\(\alpha_\ell\) alternates makes \(\eta_\ell=\alpha_\ell\beta\) alternate, which
is the disintegration clause of \cref{thm:rg-fixed-point-equations}. Finally,
\(F_0(x)=x+1\) and \(F_1(x)=x-1\) on \(\R\) have monodromy
\(F_1\circ F_0=\operatorname{id}\), which fixes every point, while neither
one-step map has a fixed point. \(\square\)

\propositionheading{Exact fixed and invariant sectors}{prop:rg-exact-fixed-sectors}
(i) The identity channel with the identity identification fixes every measure
pair. (ii) A terminal channel followed by its one-point identification has the
constant fixed theory. (iii) On \(\mathsf Y=(\R^d)^{\mathbb Z}\), let \(K_b\) be
the deterministic map
\((y_i)_i\mapsto\bigl(b^{-1/\alpha}\sum_{i\in B_j}y_i\bigr)_j\) over consecutive
blocks \(B_j\) of size \(b\), with the blocks relabeled as sites. The product of
a strictly \(\alpha\)-stable law with constant likelihood \(\mathsf m_*=M_*\varrho_*\) is a
fixed measure pair for every integer \(b\geq2\); invariance of a law for every
integer \(b\geq2\) under this normalization is the definition of strict
\(\alpha\)-stability, and at one integer \(b\) the list is not exhaustive,
because a law can be invariant for that \(b\) without being strictly stable. In
(ii) and (iii) the action is constant, whereas the identity channel, which has
every law invariant, fixes pairs with nonconstant action. (iv) Graph-flat link
data, with trivial holonomy \eqref{eq:geo-link-holonomy} on every closed walk,
stay graph flat under rooted contraction: every dressed edge
\eqref{eq:rg-dressed-boundary-edge} internal to a block is the identity, and
every coarse closed walk has trivial holonomy. (v) One-hot rows are invariant
under the softmax parameter-path derivative \eqref{eq:rg-attention-beta-softmax}.
\status{ESTABLISHED}

\noindent\emph{Proof.}
(i) and (ii) are immediate, and in (ii) the one-point space carries a constant
action. (iii) If \(Y_i\) are independent with a strictly \(\alpha\)-stable law,
\(b^{-1/\alpha}\sum_{i\in B_j}Y_i\) has the same law for every block, and the
blocks are independent, so \(\varrho_*K_b=\varrho_*\) and
\(\mathsf m_*K_b=M_*\varrho_*K_b=\mathsf m_*\). For the last clause take \(d=1\),
\(0<\alpha<1\), and the infinitely divisible law on \([0,\infty)\) with L\'evy
measure \(\nu=\sum_{k\in\mathbb Z}b^k\delta_{b^{-k/\alpha}}\) and no drift, which
exists because \(\int\min\{1,x\}\nu(dx)<\infty\)
\citep[Theorem~8.1]{Sato1999}. The image of \(\nu\) under
\(x\mapsto b^{1/\alpha}x\) is \(b\nu\), so the sum of \(b\) independent copies
has the law of \(b^{1/\alpha}\) times one copy. If the sum of \(n\) copies had
the law of \(c_n\) times one copy for every \(n\), uniqueness of the L\'evy
measure \citep[Theorem~8.1]{Sato1999} would give \(n\nu\) equal to the image of
\(\nu\) under \(x\mapsto c_nx\); matching atoms forces
\(c_n=b^{j/\alpha}\) and \(n=b^j\) for an integer \(j\), which fails for
\(n=b+1\). (iv) For an edge \(e:j\to i\) inside a block \(I\), the dressed edge
\(\Theta^x_{I\leftarrow i}\Theta^x_e(\Theta^x_{I\leftarrow j})^{-1}\) is the
holonomy of the closed walk from the root to \(j\) along the tree, across
\(e\), and back to the root, hence the identity. A product of dressed edges
along a coarse closed walk is, after the inner tree factors
\((\Theta^x_{J\leftarrow j})^{-1}\Theta^x_{J\leftarrow j'}\) are read as tree
paths inside \(J\), the holonomy of a fine closed walk, hence trivial. (v) A
row with \(\beta_J=\delta_{JK}\) has
\(\dot\beta_J=\beta_J(u_J-\sum_L\beta_Lu_L)=0\) for every \(J\). \(\square\)

Flat link data and uniform attention rows are invariant forms; they are fixed
objects only after a self-similar identification of graphs, block sizes,
receiver occupancies, and source sets is declared
(\cref{def:rg-typed-fixed-objects}).

\section{Spectra and topology dependence}
\label{sec:rg-action-spectrum-topology}

\theoremheading{Qualified positive-unital essential-spectrum theorem}{thm:rg-unital-essential-spectrum}
Let \(\mathcal V\) be a complex Banach lattice with quasi-interior unit
\(\mathbf1\in\mathcal V_+\), let \(U\in\mathcal L(\mathcal V)\) be positive and
unital, and write
\begin{equation}
r_{\mathrm{ess}}(U)=r\bigl([U]\bigr)
\label{eq:rg-calkin-essential-radius}
\end{equation}
for the spectral radius of the class \([U]\) of \(U\) in the Calkin algebra
\(\mathcal L(\mathcal V)/\mathcal K(\mathcal V)\). Assume that whenever
\(r(U)>r_{\mathrm{ess}}(U)\), the value \(r(U)\) is a pole of the resolvent and
\(U^*\) has a nonzero positive eigenfunctional \(f^*\) with
\(U^*f^*=r(U)f^*\). Then
\begin{equation}
r(U)>1\quad\Longrightarrow\quad r(U)=r_{\mathrm{ess}}(U).
\label{eq:rg-unital-essential-spectrum-result}
\end{equation}
The hypothesis \(r(U)>1\) fails for every reverse conditional expectation of
the exact arrows, which is a contraction on each \(L^p\),
\(1\leq p\leq\infty\) (\cref{thm:rg-action-lp-contraction}); the theorem
concerns positive unital operators outside the exact arrow class, and it does
not assert quasi-compactness or a Perron theorem under weaker hypotheses.
\status{ESTABLISHED}

\noindent\emph{Proof.}
The quotient map is contractive, so \(r_{\mathrm{ess}}(U)\leq r(U)\). Suppose
\(r(U)>1\) and that the inequality is strict. The hypothesis gives
\(f^*\geq0\), \(f^*\neq0\), and \(U^*f^*=r(U)f^*\). Quasi-interiority implies
\(f^*(\mathbf1)>0\): otherwise positivity makes \(f^*\) vanish on the principal
ideal generated by \(\mathbf1\), hence on its dense closure. Unitality gives
\begin{equation}
r(U)f^*(\mathbf1)
=(U^*f^*)(\mathbf1)
=f^*(U\mathbf1)=f^*(\mathbf1),
\label{eq:rg-unital-functional-pairing}
\end{equation}
a contradiction, so the two radii agree. A contraction has spectral radius
\(r(U)=\lim_n\lVert U^n\rVert^{1/n}\leq1\). \(\square\)

\propositionheading{Circle witness for norm-dependent radius}{prop:rg-circle-norm-witness}
Let \(\mathbb T=\R/\mathbb Z\) carry its geodesic metric and Haar probability
\(\mathrm{Leb}_{\mathbb T}\), let \(\mathsf T_2(z)=2z\bmod1\), and set
\begin{equation}
K(y,A)=\frac12\sum_{z\in\mathsf T_2^{-1}\{y\}}\mathbf1_A(z)
\qquad(A\subseteq\mathbb T\text{ Borel}).
\label{eq:rg-circle-forward-kernel}
\end{equation}
The kernel preserves Haar measure and its reverse is
\(K^\dagger(z,dy)=\delta_{\mathsf T_2(z)}(dy)\), so the reverse conditional
expectation is \(\mathsf Uf=f\circ\mathsf T_2\). Its spectral radius is one on
\(L^\infty(\mathrm{Leb}_{\mathbb T})\), and for every \(0<\alpha\leq1\) it is
\(2^\alpha\) on complex periodic \(C^\alpha(\mathbb T)\) with the
sup-plus-H\"older norm, which at \(\alpha=1\) is the Lipschitz space
\(C^{0,1}\). The H\"older space with this norm is not the Banach lattice of
\cref{thm:rg-unital-essential-spectrum}; the same exact kernel has different
spectral radii in two declared topologies. \status{ESTABLISHED}

\noindent\emph{Proof.}
Haar preservation gives
\(\lVert f\circ\mathsf T_2^n\rVert_\infty=\lVert f\rVert_\infty\), so
\(\lVert\mathsf U^n\rVert=1\) on \(L^\infty\). With \([f]_\alpha\) the H\"older seminorm,
\begin{equation}
[f\circ\mathsf T_2^n]_\alpha\leq2^{\alpha n}[f]_\alpha,
\qquad
\lVert\mathsf U^nf\rVert_\infty\leq\lVert f\rVert_\infty,
\label{eq:rg-circle-holder-upper-bound}
\end{equation}
so the radius on \(C^\alpha\) is at most \(2^\alpha\). For
\(f(s)=e^{2\pi is}\), the values of \(f\circ\mathsf T_2^n\) at \(0\) and
\(2^{-(n+1)}\) differ by two while their distance is \(2^{-(n+1)}\), so
\begin{equation}
[f\circ\mathsf T_2^n]_\alpha\geq2^{1+\alpha(n+1)},
\label{eq:rg-circle-holder-lower-bound}
\end{equation}
and division by the fixed \(C^\alpha\) norm of \(f\) gives the matching lower
bound. \(\square\)

\section{Two-index limits and infinite index sets}
\label{sec:rg-two-index-limits}

\definitionheading{Two-index family}{def:rg-two-index-family}
A finite-volume scale family has two indices,
\begin{equation}
X_{n,\ell},\qquad 0\leq\ell\leq L(n),
\label{eq:rg-two-index-family}
\end{equation}
where \(n\) controls system size and \(\ell\) depth. Its fixed-depth,
maximal-depth, and diagonal limits are
\begin{equation}
\lim_{\ell\to\infty}\lim_{n\to\infty}X_{n,\ell},
\qquad
\lim_{n\to\infty}X_{n,L(n)},
\qquad
\lim_{n\to\infty}X_{n,\ell(n)},
\label{eq:rg-two-index-limits}
\end{equation}
each in a declared mode of convergence. A thermodynamic scale statement
specifies a triangular family, projective consistency at fixed depth, uniform
control of the coarse error, and a diagonal scaling \(\ell(n)\); operator, law,
connection, and cross-morphism convergence are separate obligations.

\propositionheading{Fixed-depth and maximal-depth limits need not agree}{prop:rg-noncommuting-limits}
Fix an integer \(b\geq2\), put \(N=b^n\), and let \(\operatorname{Path}_N\) be the
unweighted path on \(\{1,\ldots,N\}\) with combinatorial Laplacian \(L_N\). At
depth \(0\leq\ell\leq n\) form the consecutive blocks
\begin{equation}
B_j^{(\ell)}=\{(j-1)b^\ell+1,\ldots,jb^\ell\},
\qquad1\leq j\leq b^{n-\ell};
\label{eq:rg-path-consecutive-blocks}
\end{equation}
joining two blocks exactly when a path edge crosses between them gives the
simple quotient \(\operatorname{Path}_{b^{n-\ell}}\). For \(N\) divisible by
\(b\), the unnormalized block-indicator prolongation
\begin{equation}
S_{N,b}\in\R^{N\times(N/b)},
\qquad
(S_{N,b})_{ij}=\mathbf1_{\{(j-1)b+1,\ldots,jb\}}(i),
\label{eq:rg-path-block-indicator}
\end{equation}
realizes the operator quotient
\begin{equation}
S_{N,b}^{\mathsf T}L_NS_{N,b}=L_{N/b},
\label{eq:rg-path-galerkin-quotient}
\end{equation}
and with \(S_{N,1}=I_N\) and, for \(1\leq\ell\leq n\),
\begin{equation}
S_{N,b^\ell}:=S_{N,b}S_{N/b,b}\cdots S_{N/b^{\ell-1},b},
\label{eq:rg-path-depth-block-indicator}
\end{equation}
the indicator of the blocks \eqref{eq:rg-path-consecutive-blocks},
\begin{equation}
S_{N,b^\ell}^{\mathsf T}L_NS_{N,b^\ell}=L_{N/b^\ell}.
\label{eq:rg-path-depth-galerkin-quotient}
\end{equation}
This is a precomposition of the quadratic form, not the pushforward of a law;
with orthonormalized indicators \(b^{-1/2}S_{N,b}\) the quotient is
\(b^{-1}L_{N/b}\), and at depth \(\ell\) the factor is \(b^{-\ell}\). The
eigenvalues of \(L_N\), with multiplicity, are
\begin{equation}
\lambda_{N,k}=2-2\cos\left(\frac{\pi k}{N}\right),
\qquad0\leq k\leq N-1.
\label{eq:rg-path-laplacian-eigenvalues}
\end{equation}
With the right-continuous full and zero-mode-excluded counts
\begin{equation}
\mu_N=\frac1N\sum_{k=0}^{N-1}\delta_{\lambda_{N,k}},
\qquad
\mathcal N_N(\lambda)=\mu_N((-\infty,\lambda]),
\qquad
\mathcal N^+_N(\lambda)=\frac1N
\#\{1\leq k\leq N-1:\lambda_{N,k}\leq\lambda\},
\label{eq:rg-path-spectral-counts}
\end{equation}
so that \(\mathcal N^+_N=\mathcal N_N-N^{-1}\mathbf1_{[0,\infty)}\), one has for
\(0\leq\lambda\leq4\), with
\(\mathrm k_N(\lambda)=\min\{N-1,\lfloor N\arccos(1-\lambda/2)/\pi\rfloor\}\),
\begin{equation}
\mathcal N_N(\lambda)=\frac{1+\mathrm k_N(\lambda)}{N},
\qquad
\mathcal N^+_N(\lambda)=\frac{\mathrm k_N(\lambda)}{N},
\label{eq:rg-path-exact-counts}
\end{equation}
both counts vanish for \(\lambda<0\), and for \(\lambda>4\) they equal \(1\) and
\(1-N^{-1}\). For every fixed \(\ell\), \(b^{n-\ell}\to\infty\) as
\(n\to\infty\), and both counts converge pointwise to
\begin{equation}
\mathcal N_\infty(\lambda)=\mathcal N^+_\infty(\lambda)
=
\begin{cases}
0,&\lambda<0,\\[2pt]
\displaystyle\frac1\pi\arccos\left(1-\frac{\lambda}{2}\right),
&0\leq\lambda\leq4,\\[7pt]
1,&\lambda>4,
\end{cases}
\label{eq:rg-path-thermodynamic-count}
\end{equation}
whose lower edge obeys
\begin{equation}
\mathcal N^+_\infty(\lambda)\sim\frac1\pi\sqrt{\lambda}
\qquad(\lambda\downarrow0).
\label{eq:rg-path-lower-edge-count}
\end{equation}
At maximal depth \(\ell=n\) the quotient is \(\operatorname{Path}_1\) for every
\(n\), with \(\mu_1=\delta_0\), \(\mathcal N_1=\mathbf1_{[0,\infty)}\), and
\(\mathcal N^+_1\equiv0\), so the fixed-depth thermodynamic limit and the
maximal-depth limit differ. Because \eqref{eq:rg-path-depth-galerkin-quotient}
is exact, every observable at \((n,\ell)\) depends only on \(b^{n-\ell}\), and
the diagonal \(\ell(n)=n-k\) gives \(\operatorname{Path}_{b^k}\) for every
\(n\); the example separates the limits only through this degeneracy.
\status{ESTABLISHED}

\noindent\emph{Proof.}
Each block is connected and exactly one original edge joins block \(j\) to
block \(j+1\), which proves the quotient statement. For \(z\in\R^{N/b}\),
\[
z^{\mathsf T}S_{N,b}^{\mathsf T}L_NS_{N,b}z
=\sum_{i=1}^{N-1}\bigl((S_{N,b}z)_{i+1}-(S_{N,b}z)_i\bigr)^2
=\sum_{j=1}^{N/b-1}(z_{j+1}-z_j)^2
=z^{\mathsf T}L_{N/b}z,
\]
and equality of the symmetric quadratic forms proves
\eqref{eq:rg-path-galerkin-quotient}; iteration proves
\eqref{eq:rg-path-depth-galerkin-quotient}, and scalar normalization gives the
stated factors. The discrete cosine eigenvectors of the path Laplacian give
\eqref{eq:rg-path-laplacian-eigenvalues}, and monotonicity of
\(k\mapsto\lambda_{N,k}\) gives \eqref{eq:rg-path-exact-counts}. For bounded
continuous \(f\),
\(\int fd\mu_N=N^{-1}\sum_kf(2-2\cos(\pi k/N))\to
\pi^{-1}\int_0^\pi f(2-2\cos\theta)d\theta\); the limit has the continuous
distribution function \eqref{eq:rg-path-thermodynamic-count}, and removing
the zero atom changes a count by \(N^{-1}\). Finally
\(\arccos(1-\lambda/2)\sim\sqrt\lambda\), while \(L_1=0\). The dependence on
\(b^{n-\ell}\) alone is \eqref{eq:rg-path-depth-galerkin-quotient}. \(\square\)

\propositionheading{Heat-moment susceptibility}{prop:rg-heat-susceptibility}
Let \(\mathcal N^+\) be nondecreasing and right continuous with
\(\mathcal N^+(0+)=0\), set
\[
Z_+(t)=\mathrm M_0(t),
\qquad
\mathrm M_k(t)=\int_{(0,\infty)}\lambda^ke^{-t\lambda}d\mathcal N^+(\lambda),
\]
suppose \(Z_+(t_0)<\infty\) for some \(t_0>0\), and suppose
\(\mathcal N^+(\lambda)\sim\mathrm c\lambda^\alpha\mathfrak l(1/\lambda)\) as
\(\lambda\downarrow0\), with \(\mathrm c>0\), \(\alpha>0\), and \(\mathfrak l\)
slowly varying at infinity. Then
\(\mathrm M_k(t)\sim\mathrm c\alpha\Gamma(\alpha+k)t^{-(\alpha+k)}\mathfrak l(t)\)
as \(t\to\infty\) for \(k=0,1,2\), and the entropy
\(S_+(t)=\log Z_+(t)+t\mathrm M_1(t)/Z_+(t)\) of the spectral Gibbs law
\(e^{-t\lambda}d\mathcal N^+(\lambda)/Z_+(t)\) satisfies
\begin{equation}
-\frac{dS_+}{d\log t}
=t^2\left(\frac{\mathrm M_2}{Z_+}-\left(\frac{\mathrm M_1}{Z_+}\right)^2\right)
\longrightarrow\alpha.
\label{eq:rg-heat-susceptibility}
\end{equation}
Dividing \(S_+\) by \(\log\lvert V_n\rvert\) divides this susceptibility by the
same factor, so zero-mode atoms and the entropy normalization must be declared
before a plateau is read as a spectral dimension. \status{ESTABLISHED}

\noindent\emph{Proof.}
Write \(\mathcal N^+_k(\lambda)=\int_{(0,\lambda]}\mu^kd\mathcal N^+(\mu)\).
Integration by parts and the regular variation of \(\mathcal N^+\) at zero give
\(\mathcal N^+_k(\lambda)\sim\mathrm c\alpha(\alpha+k)^{-1}
\lambda^{\alpha+k}\mathfrak l(1/\lambda)\), and the Abelian direction of the
Laplace--Stieltjes regular-variation theorem
\citep[\S1.7.1]{BinghamGoldieTeugels1987} gives
\(\mathrm M_k(t)\sim\Gamma(1+\alpha+k)\mathrm c\alpha(\alpha+k)^{-1}
t^{-(\alpha+k)}\mathfrak l(t)=\mathrm c\alpha\Gamma(\alpha+k)
t^{-(\alpha+k)}\mathfrak l(t)\). Since \(\mathrm M_1=-Z_+'\) and
\(\mathrm M_2=Z_+''\) for \(t>t_0\), differentiation gives
\(dS_+/dt=-t(\mathrm M_2/Z_+-(\mathrm M_1/Z_+)^2)\), the identity in
\eqref{eq:rg-heat-susceptibility}. The asymptotics give
\(t^2\mathrm M_2/Z_+\to\alpha(\alpha+1)\) and \(t\mathrm M_1/Z_+\to\alpha\), so the
limit is \(\alpha(\alpha+1)-\alpha^2=\alpha\). \(\square\)

This is the Abelian direction; the converse Tauberian direction is not used.

\subsection{Projective laws and refinement}
\label{sec:rg-projective-refinement}
\label{sec:cg-projective-laws}

\theoremheading{Cylinder-level projective law}{thm:cg-projective-cylinder-law}
Let \(\mathcal I\) be an index set and each coordinate space \(\mathsf Z_i\)
standard Borel. For every finite \(D\subset\mathcal I\) let \(P_D\) be a
probability law on \(\mathsf Z_D=\prod_{i\in D}\mathsf Z_i\). If
\begin{equation}
(\pi_{E,D})_\#P_E=P_D\qquad(D\subset E),
\label{eq:cg-projective-consistency}
\end{equation}
then there is a unique law \(P\) on the product cylinder sigma-algebra with
finite-dimensional marginals \(P_D\). Let a group act on
\(\mathsf Z_{\mathcal I}=\prod_{i\in\mathcal I}\mathsf Z_i\) by maps compatible
with the finite projections: either each \(g\) satisfies
\(\pi_D\circ g=g_D\circ\pi_D\) for bimeasurable bijections \(g_D\) of
\(\mathsf Z_D\), as coordinatewise product actions do, or \(g\) is an index
permutation \(\sigma\), for which
\(\pi_D\circ\sigma=\sigma_D\circ\pi_{\sigma^{-1}D}\) with the relabeling
bijection \(\sigma_D:\mathsf Z_{\sigma^{-1}D}\to\mathsf Z_D\). Then \(g_\#P\) is
the unique law with marginals \((g_D)_\#P_D\), respectively
\((\sigma_D)_\#P_{\sigma^{-1}D}\); in particular, invariance of every
finite-dimensional marginal under this rule implies \(g_\#P=P\). An action that
mixes infinitely many coordinates is not covered, because the marginals of
\(g_\#P\) are then not determined by one \(P_D\). \status{ESTABLISHED}

\noindent\emph{Proof.}
The first statement is the Kolmogorov extension theorem for standard Borel
coordinate spaces \citep{Kallenberg2021}. For a compatible action,
\((\pi_D)_\#(g_\#P)=(g_D)_\#(\pi_D)_\#P=(g_D)_\#P_D\) for every finite \(D\),
and the permutation case is the same computation with
\(\pi_D\circ\sigma=\sigma_D\circ\pi_{\sigma^{-1}D}\); uniqueness of the
extension identifies \(g_\#P\). \(\square\)

Consistency alone does not put the limit on smooth sections, define a
continuum reference density, justify an ELBO limit, or produce a scale flow
(\cref{open:prob-continuum-theory}), and an energy recursion does not supply
compatible laws: a probabilistic scale flow needs normalized kernels with
\(\mu_k=\mu_\ell K_{\ell k}\) and \(K_{\ell k}K_{kr}=K_{\ell r}\), with
corresponding rules for the two bundle sectors and their reference measures.
For a countable agent or time index the extension supplies the normalized
joint, and the finite identities of the main manuscript hold for every regular
finite-cylinder observation and finite block. Conditioning on an infinite
record needs a chosen regular conditional at an almost-everywhere record or a
proved DLR, Doob-transform, or finite-volume posterior limit, and an
infinite-law relative-entropy pushforward needs absolute continuity of the
recognition law with respect to the posterior on the infinite product and an
infinite coarse channel; \cref{ch:constructive-spatial-rg,ch:geometry-pressure}
prove such hypotheses on their own domains, whose generality is
\cref{open:rg-spatial-generality}.

\propositionheading{Infinite records and refinement}{prop:rg-infinite-records-refinement}
(i) A projectively consistent family can give a prescribed infinite record
probability zero: for independent binary successes of probability
\(p\in(0,1)\) each, the record of all successes has probability zero. (ii) If
a coarse map \(Q\) is noninjective, no refinement kernel \(R\) recovers every
realized microstate, and a refinement kernel compatible with \(Q\) is a
stochastic right section,
\begin{equation}
Q_\#R(c,\cdot)=\delta_c,
\label{eq:rg-stochastic-right-section}
\end{equation}
which redraws a compatible microstate. \status{ESTABLISHED}

\noindent\emph{Proof.}
(i) The cylinder events of the first \(n\) successes decrease to the full
record and have probabilities \(p^n\to0\). (ii) If \(Q(x)=Q(x')\) with
\(x\neq x'\), recovery of both would require
\(R(Q(x),\cdot)=\delta_x\) and \(R(Q(x),\cdot)=\delta_{x'}\). \(\square\)

\propositionheading{A conditional operational diagonal}{prop:rg-refresh-operational-diagonal}
Consider targets $P^\star_N$ on finite networks at $c_*$, evaluated at depths
$\ell(N)\to\infty$. Suppose each admits the fixed-target estimate of
\cref{prop:rg-conditional-refresh-entropy}, with $m_N$ blocks, refresh
floor $\mathrm a_N>0$, entropy factorization constant $\mathrm C_N\geq1$,
and finite initial entropy $e_{0,N}$. Take $k_N$ inference updates and
numerical kernel errors $\varepsilon_{j,N}$ as in
\cref{prop:rg-refresh-target-transfer}. A sufficient condition for vanishing
ideal endpoint entropy and numerical-to-target TV is
\begin{equation}
\frac{k_N\mathrm a_N}{m_N\mathrm C_N}-\log(1+e_{0,N})\longrightarrow+\infty,
\qquad
\sum_{j<k_N}\varepsilon_{j,N}\longrightarrow0.
\label{eq:rg-refresh-operational-diagonal}
\end{equation}
If, in addition, specified readout channels $C_N^{\mathrm{out}}$ take these laws to one
Polish comparison space and $P^\star_NC_N^{\mathrm{out}}$ converges weakly to a law
$P^\star_\infty$, then the numerical readout laws have the same weak limit.
This is a transfer conditional on target convergence and on the displayed
uniform estimates; it supplies neither of those premises, surviving
parents, nor an operational cost bound.
\status{ESTABLISHED}

\noindent\emph{Proof.}
The fixed-target entropy recurrence gives
$e_{k_N,N}\leq e_{0,N}\exp[-k_N\mathrm a_N/(m_N\mathrm C_N)]$,
which tends to zero under the first condition. Pinsker's inequality and
\eqref{eq:rg-refresh-numerical-transfer} give vanishing numerical TV.
Readout contracts TV, so the expectations of every bounded continuous
function differ by at most its oscillation times that TV distance.
Target weak convergence therefore transfers to the numerical readout.
\hfill$\square$

\openproblemheading{Two-index limits and universality}{open:rg-two-index-limits}
Fixed-depth and maximal-depth limits can differ
(\cref{prop:rg-noncommuting-limits}), finite networks at \(c_*\) do not inhabit
the asymptotic classification (\cref{prop:rg-finite-depth-no-asymptotics}),
and the lattice part fixes every finite blocking word before its volume limit
on its own domains (\cref{ch:geometry-pressure,open:rg-spatial-generality}). It
is not known whether a triangular family of finite agent networks at \(c_*\)
has fixed-depth or diagonal limits with nonvacuous comparison maps and
nontrivial exponents, whether the two limits can be exchanged or a nontrivial
diagonal limit exists, or whether the spectral labels of
\cref{def:rg-gaussian-spectral-labels} converge for a declared Gaussian
hierarchy. A universality statement needs scheme-invariant observables
(\cref{prop:rg-scheme-dependence}) with proved convergence; one matched
exponent does not give it, since the limiting measure, amplitudes, corrections
to scaling, invariant faces, basin, and blocking scheme may still differ. The
problem is settled by a triangular family with projective consistency at fixed
depth, uniform control of the coarse error, a diagonal scaling \(\ell(n)\), and
a proof of convergence of a declared invariant observable, or by an
impossibility theorem for a stated class. The conditional operational
diagonal of \cref{prop:rg-refresh-operational-diagonal} transfers an
already established target limit; it does not construct the limiting
nonlinear law or its compatible retained dynamics.

\section{Gaussian operator flow on infinite hierarchies}

This section applies the preceding types to the Gaussian interaction family
of \cref{def:gauss-interaction-family}, with the hard-aggregation map
\eqref{eq:grg-aggregation} and the positive rescaling \eqref{eq:grg-flow}, whose
finite exact properties are established in
\cref{ch:gaussian-renormalization}. Along a hierarchy the source and target
dimensions differ, so \eqref{eq:grg-flow} is a cocycle of positive maps in the
sense of \cref{def:rg-dynamical-type}; a fixed ray is typed only after
comparison maps identify a repeated hierarchy with one parameter space.

\subsection{Dense scaling and the homogeneous sector spectrum}
\label{sec:grg-numerical-gate}

\propositionheading{Dense scaling and topology-dependent exponents}{prop:rg-gaussian-dense-scaling}
(i) On a homogeneous complete graph with equal blocks of size \(b\), common
self term \(A\), and common coupling \(W\), the unrescaled aggregation acts on
the pair \((A,W)\) by
\begin{equation}
A\longmapsto bA,\qquad W\longmapsto b^2W.
\label{eq:grg-homogeneous-map}
\end{equation}
On \(\Sym^K\oplus\Sym^K\) its eigenvalues are \(b\) on the self sector and
\(b^2\) on the coupling sector, each of multiplicity \(K(K+1)/2\), so the
spectral ratio outside the dominant eigenspace is
\begin{equation}
\frac{b}{b^2}=\frac1b.
\label{eq:grg-sector-gap}
\end{equation}
(ii) The dense normalization \(\zeta=b^2\) in \eqref{eq:grg-flow} fixes the
coupling sector and contracts the self sector by \(1/b\), so the candidate
limiting ray has \(A=0\). It is a boundary operator: the consensus directions
lie in its kernel, so it defines no normalized Gaussian density on the full
space. (iii) For a typical block \(B(b)\) and coarse cut \(\partial B(b)\), let
\(W_\partial(b)=\sum_{e\in\partial B(b)}W_e\) and
\(A_B(b)=\sum_{i\in B(b)}A_i\). If
\(\lVert W_\partial(b)\rVert\asymp b^{\xi_W}\) and
\(\lVert A_B(b)\rVert\asymp b^{\xi_A}\) in a declared matrix norm, a scalar
normalization that keeps the coupling sector nondegenerate scales as
\(b^{\xi_W}\), and the normalized self sector as \(b^{\xi_A-\xi_W}\). If
moreover the cut has \(n_{\mathrm{cut}}(b)\asymp b^\xi\) edges,
\(W_\partial(b)/n_{\mathrm{cut}}(b)\to\overline W\neq0\), and
\(A_B(b)/b\to\overline A\neq0\), then \(\xi_W=\xi\) and \(\xi_A=1\), so the
coupling and normalized self exponents are \(\xi\) and \(1-\xi\); a cut between
two blocks of the complete graph has \(\xi=2\), and consecutive blocks of a
nearest-neighbor chain have \(\xi=0\). \status{ESTABLISHED}

\noindent\emph{Proof.}
(i) By \eqref{eq:grg-aggregation}, a coarse block \(I\) has self term
\(\sum_{i\in I}A_i=bA\), and two coarse blocks \(I\neq J\) have coupling
\(\sum_{i\in I,j\in J}W_{ij}=b^2W\); the linear map is diagonal on the two
sectors. (ii) Under \(\zeta=b^2\) the self term scales by \(b/b^2\). With
\(A=0\), the block row \(i\) of the assembled precision applied to
\(\mathbf1\otimes v\) is
\((A_i+\sum_{j\neq i}W_{ij})v-\sum_{j\neq i}W_{ij}v\), which equals
\(A_iv=0\), so every
consensus vector lies in the kernel and \(\exp(-x^{\mathsf T}\Lambda x/2)\) is
not integrable along those directions. (iii) With \(\zeta(b)\asymp
\lVert W_\partial(b)\rVert\), the normalized sectors are of orders \(1\) and
\(b^{\xi_A-\xi_W}\); the count assumptions give
\(\lVert W_\partial(b)\rVert\sim n_{\mathrm{cut}}(b)\lVert\overline W\rVert\asymp
b^\xi\) and \(\lVert A_B(b)\rVert\sim b\lVert\overline A\rVert\). A cut between
two blocks of size \(b\) of the complete graph has \(b^2\) edges, and a chain
cut has one. \(\square\)

Pinning, quotienting the kernel, and retaining an external positive mass are
different repairs of the boundary ray and define different probability
problems; inserting a pseudoinverse into the original Lebesgue density is not
a further repair. Edge counts alone do not determine the exponents when the
fine weights vary with \(b\), and relevance of the self term is a property of
the declared weighted graph family and normalization, not of
\eqref{eq:grg-flow} alone.

\subsection{Cone dynamics}
\label{sec:grg-cones}

Let
\begin{equation}
\mathcal W_N
=\left(\prod_i\PSD^K\right)\times\left(\prod_{i<j}\PSD^K\right)
\label{eq:grg-full-cone}
\end{equation}
be the full parameter cone. On the interior of a proper cone \(\mathcal W\),
Hilbert's projective metric is
\begin{equation}
d_H(x,y)=\log\frac{t_+(x,y)}{t_-(x,y)},
\qquad
t_+=\inf\{t>0:x\preceq_{\mathcal W}ty\},
\quad
t_-=\sup\{t>0:ty\preceq_{\mathcal W}x\}.
\label{eq:grg-hilbert}
\end{equation}
A positive endomorphism of finite projective diameter contracts this metric; a
matrix with strictly positive entries on the nonnegative orthant is the
familiar case, and for a primitive nonnegative matrix the contraction applies
to a strictly positive power, since the matrix itself can map boundary rays to
boundary rays and then has infinite projective diameter
\citep{Birkhoff1957,nussbaum1988hilbert}. These hypotheses fail for sum-only
aggregation on the full cone.

\propositionheading{The full-cone route is blocked}{prop:grg-full-cone-blocked}
After a hierarchical identification makes \(\mathcal A_S\) an endomorphism,
sum-only aggregation is not primitive on \(\mathcal W_N\), and the dense
candidate ray with \(A=0\) of \cref{prop:rg-gaussian-dense-scaling} lies on
\(\partial\mathcal W_N\), whereas a full-cone Birkhoff contraction has an
interior Perron ray; Birkhoff's full-cone theorem therefore cannot prove
attraction to that candidate. On the coupling cone alone, if an identified
spatial map \(B\geq0\) has a positive eigenvector \(c=(c_e)_e\), then for every
\(M\in\PSD^K\)
\begin{equation}
W_e=c_eM
\label{eq:grg-inert-fiber}
\end{equation}
is an eigenray with the same spatial eigenvalue, so sum-only aggregation does
not mix matrix rays and cannot select a fiber direction. \status{ESTABLISHED}

\noindent\emph{Proof.}
An absent coarse cut remains a zero edge coordinate, so the map can meet the
boundary even from positive data on the represented fine edges. For a subspace
\(U\subseteq\R^K\), the face
\begin{equation}
\mathcal F_U=\{(W_e)_e:\operatorname{ran}W_e\subseteq U\ \text{for every }e\}
\label{eq:grg-common-range-face}
\end{equation}
is invariant, because the range of a sum of positive semidefinite matrices
with ranges in \(U\) lies in \(U\), and proper invariant faces rule out
primitivity. The coordinate \(A=0\) is a boundary coordinate of
\(\prod_i\PSD^K\), while the Perron ray of a strict interior contraction is
interior. Finally, aggregation adds weights, so
\((B\otimes\operatorname{id})(c\otimes M)=(Bc)\otimes M\). \(\square\)

\propositionheading{Fixed-\(B\) attraction under a primitive identification}{conj:grg-fixed-b-attraction}
Suppose a hierarchical identification makes one aggregation step act on
coupling coordinates as \(B\otimes\operatorname{id}\), that is
\(W'_f=\sum_eB_{fe}W_e\), for one fixed nonnegative matrix \(B\) on edge indices
with \(B^q\) strictly positive for some \(q\geq1\). Let \(r_B>0\) be its spectral
radius and \(u,v>0\) its left and right Perron vectors with \(u^\top v=1\).
(i) Fix \(M_0\succ0\). On the scalarized cone
\begin{equation}
\mathcal W_{M_0}=\{(c_eM_0)_e:c_e\geq0\},
\label{eq:grg-scalarized-cone}
\end{equation}
every nonzero orbit converges projectively to the ray \((v_eM_0)_e\), which is
therefore the unique attracting ray; if \(v_{ij}=x_ix_j\), that ray is the
biadditive coupling \(W_{ij}=x_ix_jM_0\) of \eqref{eq:grg-biadditive}.
(ii) On the full coupling cone, for every initial datum \((W_e^{(0)})_e\) with
\(W_e^{(0)}\succeq0\),
\begin{equation}
r_B^{-\ell}\bigl(B^\ell W^{(0)}\bigr)_f\longrightarrow v_fM_\infty,
\qquad
M_\infty=\sum_eu_eW_e^{(0)},
\label{eq:grg-full-cone-perron-limit}
\end{equation}
and \(\operatorname{ran}M_\infty=\sum_e\operatorname{ran}W_e^{(0)}\). Every
nonzero orbit converges projectively to the product ray through
\(v\otimes M_\infty\): aggregation fixes the fiber matrix as the left-Perron
average of the initial weights and does not select it. \status{ESTABLISHED}

\noindent\emph{Proof.}
A nonnegative matrix with a strictly positive power is primitive, and then
\(r_B^{-\ell}B^\ell\to vu^\top\) \citep[\S8.5]{HornJohnson2013}. Hence
\(r_B^{-\ell}(B^\ell W^{(0)})_f\to v_f\sum_eu_eW_e^{(0)}\). Because every
\(u_e>0\) and every \(W_e^{(0)}\succeq0\),
\(\ker M_\infty=\bigcap_e\ker W_e^{(0)}\), which gives the range identity and
\(M_\infty\neq0\) for every nonzero datum, so the normalized orbit converges to
the product ray. On \(\mathcal W_{M_0}\) the limit reads \(v(u^\top c)M_0\), with
\(u^\top c>0\) for every nonzero \(c\geq0\). Birkhoff's contraction applied to the
strictly positive matrix \(B^q\) gives the scalar conclusion with a geometric
rate in Hilbert's metric \citep{Birkhoff1957}. \(\square\)

The proposition covers one fixed identification acting on edge indices. It
says nothing about a nonautonomous sequence \(B_\ell\), whose projective
contraction can approach a periodic orbit
(\cref{prop:rg-contraction-no-fixed-ray}), or about an identification that
also acts on fiber directions by a level-dependent congruence, for which the
step is not \(B\otimes\operatorname{id}\).

\openproblemheading{Primitive hierarchical identification}{open:rg-scalarized-attraction}
\Cref{conj:grg-fixed-b-attraction} proves attraction for one fixed primitive
identification \(B\otimes\operatorname{id}\). It is not known whether sum-only
aggregation on a declared family of graphs admits a hierarchical
identification that yields one fixed primitive \(B\) and, for the biadditive
reading, one whose Perron vector factorizes as \(v_{ij}=x_ix_j\); attraction,
uniqueness, and scheme independence of the dense ray of
\cref{prop:rg-gaussian-dense-scaling} likewise need a declared invariant cone,
hierarchy, and limiting protocol. The problem is settled by constructing such
an identification for a stated family of graphs, or by proving that none
exists.

\subsection{Thermodynamic spectral labels}

\definitionheading{Thermodynamic spectral labels}{def:rg-gaussian-spectral-labels}
For a declared sequence \((V_n,L_n,R_n)\) of interaction Laplacians and
positive reference forms, with generalized eigenvalues \(\lambda^{(n)}_k\) of
the pencil \((L_n,R_n)\) of \cref{sec:grg-spectra}, the integrated density of
states is
\begin{equation}
\mathcal N_n(t)
=\frac{1}{\lvert V_n\rvert K}
\#\{k:\lambda_k^{(n)}\leq t\}.
\label{eq:grg-ids}
\end{equation}
A cumulative exponent, \(\mathcal N(t)-\mathcal N(0)\sim\mathrm ct^\alpha\), is a
label of the sequence only after convergence \(\mathcal N_n\to\mathcal N\) in a
stated mode has been proved.

A fixed finite graph has an atomic spectral measure, so no density exponent is
defined there. Depth and size form the two-index family of
\cref{def:rg-two-index-family}, and neither the existence nor the commutation
of its limits follows from the finite congruence identity
(\cref{open:rg-two-index-limits}).
